\section{上线之后系统开始改变它所预测的世界}
\label{sec:17-Synthesis/01-System-Lifecycle/06}

新闻推荐系统全量上线的第一周，人均点击涨了 \(8\%\)。第六周开始回落，第十周低于上线前。同一时期，编辑团队注意到一件事：某一类耸动标题的内容在曝光里的占比从 \(12\%\) 涨到了 \(47\%\)。

没有人改过模型。每周的重训用的都是同一套代码、同一套超参，唯一变化的是训练数据——而训练数据现在来自上一周由这个模型决定的曝光。用户只能点他看得见的东西，模型只能从他点了什么来学，于是模型每一轮都在确认自己上一轮的判断。十周之后它对世界的看法已经与真实的用户偏好没有多少关系了。

风控那一侧有一个结构完全相同、症状相反的例子。被模型拒掉的申请不会放款，因而永远不会产生还款结果；下一轮重训时，那片区域里一个正样本也没有，模型只会更确信自己拒对了。它的置信度稳步上升，而它对被拒人群的判断从来没有被检验过。

前面五篇都建立在一个未说出口的前提上：世界在那里，系统去猜它。这个前提在上线的那一刻失效。\textbf{部署把一个预测问题变成了一个反馈系统}，而反馈系统的行为不能用预测的语言分析。

\begin{center}
\begin{tikzpicture}[x=1cm,y=1cm,font=\small,>=stealth]
  \draw[black!45,fill=blue!8] (0.5,2.8) rectangle (3.3,3.8);
  \node[align=center] at (1.9,3.3) {模型};
  \draw[black!45,fill=blue!8] (5.2,2.8) rectangle (8.4,3.8);
  \node[align=center] at (6.8,3.3) {决策与干预};
  \draw[black!45,fill=orange!10] (5.2,0.4) rectangle (8.4,1.4);
  \node[align=center] at (6.8,0.9) {环境的响应};
  \draw[black!45,fill=orange!10] (0.5,0.4) rectangle (3.3,1.4);
  \node[align=center] at (1.9,0.9) {新收集的数据};
  \draw[->,black!60] (3.3,3.3) -- (5.2,3.3);
  \draw[->,black!60] (6.8,2.8) -- (6.8,1.4);
  \draw[->,black!60] (5.2,0.9) -- (3.3,0.9);
  \draw[->,red!70!black,very thick] (0.5,0.9) -- (-0.25,0.9) -- (-0.25,3.3) -- (0.45,3.3);
  \node[rotate=90,anchor=south,font=\footnotesize,red!70!black] at (-0.4,2.1) {分布被自己改写};
  \node[align=left,anchor=west,font=\footnotesize,black!70] at (8.9,3.3)
    {阈值与代价决定\\这里输出什么（第一篇）};
  \node[align=left,anchor=west,font=\footnotesize,black!70] at (8.9,2.1)
    {这一支箭是干预，\\不是观察};
  \node[align=left,anchor=west,font=\footnotesize,black!70] at (8.9,0.9)
    {被拒绝的那些\\永远不产生标签};
\end{tikzpicture}

\emph{前五篇都默认那条红色的回边不存在。上线把它接上了，这张图于是不再是一条流水线，而是一个会自己演化的动力系统。}
\end{center}

\subsection{三种漂移，三种诊断信号}

先把不需要反馈就能发生的那部分讲清楚，因为它是基线。把输入记作 \(X\)、标签记作 \(Y\)，训练时的联合分布是 \(P_{\text{train}}(X,Y)\)，线上是 \(P_{\text{live}}(X,Y)\)。这两个分布可以在三个不同的地方分家。

边缘分布 \(P(X)\) 变了而条件 \(P(Y\given X)\) 没变，是协变量漂移：用户构成变了、新渠道带来了另一类人，但``给定这些特征，违约概率是多少''这条规律还成立。它的诊断信号最友好——只看输入就能发现，不需要等标签回来。逐特征比较训练与线上的分布，或者训一个分类器去区分两批输入（能区分得开就说明漂了），都是常用做法。

\(P(Y)\) 变了是标签漂移：欺诈的整体比例季节性上升，而每类样本长什么样没变。它需要标签才能诊断，而标签通常滞后数周甚至数月。

最难的是概念漂移：\(P(Y\given X)\) 本身变了，同样的特征对应了不同的结果。诈骗团伙换了打法、政策调整改变了还款能力，这时输入分布可以纹丝不动，而模型学到的映射已经过期。它在输入侧完全看不出来，只能靠持续回收的标签去发现，而且发现时损失已经产生了。这三者需要的应对不同，混为一谈会导致做了一堆输入监控却挡不住真正的失效。

要判断``分布在变''这句话本身是否有意义，得先有平稳性的概念。\hyperref[sec:06-Random-Processes/01-Foundations/03]{``均值、相关与平稳性''}给出的区分是：弱平稳只要求均值与自相关不随时间移动，而一段观测看起来平稳既可能是真的平稳，也可能是观测窗口太短。所以监控里那条``均值没变''不足以支撑``没有漂移''，尤其在存在缓慢趋势时。

\subsection{模型自己造成的漂移不是漂移}

上面三种都可以在没有部署的情况下发生——世界自己在变。反馈回路造成的是另一回事：分布之所以变，正是因为模型在按它自己的预测行动。

这时``预测''与``干预''的界限消失了。\hyperref[sec:13-Machine-Learning/12-Causal-Inference/01]{``观察、干预与反事实回答不同问题''}把这三层严格分开：观察问的是``看到 \(X=x\) 的那些个体，\(Y\) 是多少''，干预问的是``如果把 \(X\) 设成 \(x\)，\(Y\) 会是多少''，反事实问的是``对这个已经发生了别的事的个体，如果当初设成 \(x\) 会怎样''。一个上线的系统做的是第二件事，而它的训练与评估全都建立在第一件事上。

这不是术语讲究。同一个联合分布可以对应多张不同的因果图，而观察数据无法在它们之间区分——\hyperref[sec:13-Machine-Learning/12-Causal-Inference/02]{``结构因果模型把干预写成机制替换''}说明了差别在哪里：干预是把某个变量的生成机制整个换掉（\(\operatorname{do}(X=x)\)），这个操作在联合分布里根本没有对应的写法，必须先有图才能算。所以``模型预测得准''与``按模型行动会得到想要的结果''之间没有推理关系，两者需要的假设不是同一组。

具体到该控制什么变量，\hyperref[sec:13-Machine-Learning/12-Causal-Inference/03]{``混杂、碰撞点与后门调整''}给出的判断在任何回归解释里都用得上：控制混杂消除偏差，控制碰撞点\emph{制造}偏差。多控制几个变量并不更保险——把一个共同结果放进回归，会在两个原本独立的原因之间凭空造出关联。这条错误在特征工程阶段几乎不可能被指标发现，因为它并不损害预测精度，只损害解释和干预。

即使因果效应在数学上可识别，也还有最后一层。\hyperref[sec:13-Machine-Learning/12-Causal-Inference/06]{``可识别不等于可外推也不等于可决策''}把三件事分开：能不能从这份数据算出效应、这个效应能不能搬到另一群人身上、以及知道了效应之后按它行动是否划算。推荐系统上线时，前一个人群上估出来的效应正在被系统自己改变的人群所取代，第二层就已经站不住了。

\subsection{把自己的输出喂回输入的系统会不会震荡}

一旦承认那条红边存在，问题的性质就变了：这不再是统计问题，而是稳定性问题。而稳定性正是控制论处理了一百年的东西。

\hyperref[sec:16-Control-and-Reinforcement-Learning/01-Optimal-Control/01]{``从开环到闭环反馈''}给出的第一条判断是：开环策略把整条动作序列事先算好，闭环策略根据观测到的状态随时修正，两者面对扰动时的表现天差地别。一个每周重训的模型是闭环的，只是它的``反馈''慢到以周计，而且没有人为它设计过增益。反馈系统的经典教训在这里原样适用——反馈太强会震荡，延迟太大会让原本稳定的回路失稳，而这两件事都与模型精度无关。

\hyperref[sec:15-Differential-Equations-and-Dynamical-Systems/02-Stability-and-Dynamical-Systems/02]{``稳定性与 Lyapunov 函数''}提供的是分析工具：找一个沿着系统演化单调下降的标量函数，就能不解方程地断言轨迹会收敛到哪里。这在部署监控里可以退化成一个很朴素的做法——挑一个应当保持有界的量（曝光的内容多样性、放款人群的特征分布与全体申请人群的距离、策略与上一版的差异），持续盯着它是不是单调恶化。它给不出证明，但它把``系统在往哪里去''变成了一个可观测的问题。

\hyperref[sec:15-Differential-Equations-and-Dynamical-Systems/02-Stability-and-Dynamical-Systems/03]{``离散动力系统、分岔与混沌''}补上另一半：每周重训是一个离散时间迭代，而离散迭代的稳定条件比连续系统更苛刻，参数缓慢变化时系统可以在某一刻突然改变定性行为。上线后``指标先涨后跌''的曲线形状，与分岔前后的行为切换是同一类现象——不是某一天出了故障，是回路增益越过了某个临界值。

强化学习那一侧早就在正面处理这些问题，它的经验对任何闭环系统都成立。\hyperref[sec:16-Control-and-Reinforcement-Learning/02-Reinforcement-Learning/08]{``Off-policy 校正用权重交换分布偏差与方差''}讲的是怎样用旧策略收集的数据评估新策略：重要性权重在原理上给出无偏估计，代价是两个策略差别一大，权重的方差就爆炸。这正是``用上一版模型的日志评估下一版''这件事的数学形态，也解释了为什么日志评估在策略改动大时不可信。\hyperref[sec:16-Control-and-Reinforcement-Learning/02-Reinforcement-Learning/06]{``Deadly Triad 让自举误差在分布外放大''}给出的则是一条更尖锐的警告：函数逼近、自举、离策略数据这三样凑齐时，误差会在分布外区域自我放大而不是被平均掉。上面那个风控例子就是这个结构的最简版本——模型在自己没有数据的区域自我确认。

\hyperref[sec:06-Random-Processes/03-Markov-Chains/04]{``平稳分布、可逆性与遍历收敛''}给了这一切一个乐观的读法：一个不断被自己改写的系统，如果转移规则稳定，长期确实会稳定到某个分布上。问题只是没有任何东西保证那个分布是我们想要的——它可能就是全是耸动标题的那个分布。系统会收敛，但收敛到哪里由回路的结构决定，与初始的模型质量关系不大。

\subsection{这一步做错时，故障长什么样}

模型上线后指标先涨后跌是最典型的形态。前几周的涨幅是真的——那时训练数据还来自旧策略，模型确实在一个未被污染的分布上学到了东西。跌是回路开始生效的信号。看到这条曲线时，先查内容或人群构成的集中度，而不是先怀疑模型退化。

重训之后越来越极端：每一轮曝光都更集中、每一轮拒绝都更严格、每一轮推荐都更同质。这是自我确认回路，特征是模型的置信度上升而实际表现下降。破它需要主动往回路里注入不受模型控制的数据——推荐里的随机探索位、风控里的小额随机放款、排序里的固定比例流量。这笔成本是买回那部分被系统自己抹掉的分布，性质上与买保险相同。

两个模型互相喂数据导致同质化：定价系统读取竞品价格、竞品读取你的价格，或者内容生成模型的输出成为下一代的训练语料。这类回路的时间常数往往很长，等症状明显时已经不可逆，因为多样性一旦从数据里消失，就没有地方再把它找回来。

还有一种最容易被误判为成功：所有监控指标平稳，模型置信度高，投诉率低——因为系统已经把可能产生分歧的那部分输入全部排除在外了。它不是稳定，是收敛到了一个自洽而狭窄的不动点。判断它的办法只有一个，就是看还有多少数据是这个系统没有参与决定的。

\subsection{这条生命周期给出的是一张对账表}

走完这六站，可以回头看整条路径真正说了什么。

诉求变成目标函数时丢掉的东西，训练再好也补不回来；数据里被采样机制和预处理悄悄用掉的假设，不会在任何一个指标上报警；模型族的假设写错了不会报错，只会让结论有偏；训练只保证在那个目标函数下找到了好点，不保证目标写对了；评估背书的是一句关于未见数据的断言，而背书的对象是程序不是那一次的结果；上线之后，前五步共同依赖的``世界在那里''这个前提本身失效了。

每一站的故障都在下游变形之后才被看见，这就是为什么现实中排查总要逆着这条链往回走。而逆着走的能力，恰恰来自知道每一站承诺了什么、没承诺什么——这本书前十六个部分讲的就是这些承诺各自的成立条件。

这一路是按时间顺序切的：一个系统从诞生到改变世界，依次经过哪些环节。还有另一种切法，把同一批工具横过来看——会发现``用局部信息控制全局''这件事在极限、凸性、Lyapunov 函数、Bellman 方程里是同一个想法的四次出现；``把复杂对象投影到低维再看''在最小二乘、PCA、变分推断、模型压缩里也是同一个动作。下一个单元做的就是这件事：不按流程，按反复出现的少数几个想法重新组织全书，看看它们各自在多少个学科里换过名字。
